% ======================================================================
\section{Linear recurrences and powers of \texorpdfstring{$2\times2$}{2x2} matrices}\label{app:linear}
% ======================================================================

These two standard facts are used in \cref{prop:noise-floor,prop:quadratic-stability}.

\begin{lemma}[Second-order linear recurrences]\label{lem:recurrence}
Let $a\in\R$, $b\ne0$, and let $\lambda_1,\lambda_2\in\mathbb C$ be the roots of
$p(\lambda)=\lambda^2-a\lambda+b$. The solutions $(x_k)_{k\ge-1}$ of
$x_{k+1}=ax_k-bx_{k-1}$ ($k\ge0$) are exactly the sequences
\begin{enumerate}[label=\textup{(\roman*)}]
\item $x_k=c_1\lambda_1^k+c_2\lambda_2^k$ with $c_1,c_2\in\mathbb C$, if $\lambda_1\ne\lambda_2$;
\item $x_k=(c_1+c_2k)\lambda_1^k$ with $c_1,c_2\in\mathbb C$, if $\lambda_1=\lambda_2$.
\end{enumerate}
\end{lemma}

\begin{proof}
\emph{These are solutions.} For a root $\lambda$,
$\lambda^{k+1}-a\lambda^k+b\lambda^{k-1}=\lambda^{k-1}p(\lambda)=0$, and sums of solutions
are solutions. In case (ii), $\lambda_1=a/2$ and $b=\lambda_1^2$ (product of the roots), and
\[
  (k+1)\lambda_1^{k+1}-ak\lambda_1^k+b(k-1)\lambda_1^{k-1}
  =\lambda_1^{k-1}\bigl[k\,p(\lambda_1)+\lambda_1^2-b\bigr]=0 .
\]
\emph{There are no others.} A solution is determined by $(x_{-1},x_0)$, so it suffices to
show that every pair $(x_{-1},x_0)$ is attained. In case (i) we need
$c_1\lambda_1^{-1}+c_2\lambda_2^{-1}=x_{-1}$, $c_1+c_2=x_0$; the determinant of this linear
system is $\lambda_2^{-1}-\lambda_1^{-1}=\frac{\lambda_1-\lambda_2}{\lambda_1\lambda_2}\ne0$
(the roots are nonzero because $\lambda_1\lambda_2=b\ne0$). In case (ii) we need
$(c_1-c_2)\lambda_1^{-1}=x_{-1}$ and $c_1=x_0$, which has the solution $c_1=x_0$,
$c_2=x_0-\lambda_1x_{-1}$.
\end{proof}

When $b=\beta=0$ the heavy-ball recurrence reduces to $x_{k+1}=ax_k$, with solutions
$x_k=a^kx_0$; the arguments of \cref{prop:quadratic-stability} apply verbatim with the roots
$a$ and $0$.

\begin{lemma}[Powers of a $2\times2$ matrix]\label{lem:matrix-powers}
Let $M$ be a real $2\times2$ matrix whose eigenvalues satisfy
$\max\{|\lambda_1|,|\lambda_2|\}<\rho<1$. Then there is a constant $C_M$ such that every entry
of $M^k$ is at most $C_M\rho^k$ in absolute value, for all $k\ge0$. Consequently, for any
$2\times2$ matrices $C_0,Q$, we have $M^kC_0(M^\top)^k\to0$, and the series
$\sum_{j\ge0}M^jQ(M^\top)^j$ converges.
\end{lemma}

\begin{proof}
If $\lambda_1\ne\lambda_2$, $M$ is diagonalizable over $\mathbb C$: $M=V\,\mathrm{diag}(\lambda_1,\lambda_2)V^{-1}$,
so $M^k=V\,\mathrm{diag}(\lambda_1^k,\lambda_2^k)V^{-1}$ and its entries are bounded by a constant
times $\max_i|\lambda_i|^k\le\rho^k$. If $\lambda_1=\lambda_2=\lambda$, then either
$M=\lambda I$, or $M=VJV^{-1}$ with the Jordan block $J=\begin{pmatrix}\lambda&1\\0&\lambda\end{pmatrix}$,
for which $J^k=\begin{pmatrix}\lambda^k&k\lambda^{k-1}\\0&\lambda^k\end{pmatrix}$; since
$k|\lambda|^{k-1}/\rho^k\to0$, the entries of $M^k$ are again bounded by a constant times
$\rho^k$. The entries of $M^jQ(M^\top)^j$ are then bounded by a constant times $\rho^{2j}$, so
the series converges absolutely, and similarly $M^kC_0(M^\top)^k\to0$.
\end{proof}

% ======================================================================
\section{Reproducing the figures}\label{app:reproduce}
% ======================================================================

The repository contains the \LaTeX{} source of this report (\texttt{report/}) and the
Python code for all figures (\texttt{experiments/}, requiring \texttt{numpy} and
\texttt{matplotlib}). From the repository root:
\begin{verbatim}
  pip install -r experiments/requirements.txt
  python experiments/run_all.py          # regenerates report/figures/*.pdf
  cd report && latexmk -pdf main.tex     # rebuilds the PDF
\end{verbatim}
or simply \texttt{make}. The scripts and the figures they produce are:
\begin{center}\small
\begin{tabular}{@{}lll@{}}
\toprule
script & figure & content\\
\midrule
\texttt{potential.py} & \cref{fig:potential} & $f(x_k)$ versus $\Phi_k$; checks monotonicity of $\Phi_k$\\
\texttt{regions.py} & \cref{fig:regions} & parameter regions\\
\texttt{sgd\_vs\_sgdm.py} & \cref{fig:sgd-vs-sgdm} & SGD versus SGDM on (P3)\\
\texttt{noise\_floor.py} & \cref{fig:noise-floor} & noise floor versus $\gamma$\\
\texttt{diminishing.py} & \cref{fig:diminishing} & diminishing step sizes\\
\texttt{counterexample.py} & \cref{fig:lrp} & the $3$-cycle of \cref{ex:lrp}, verified exactly\\
\bottomrule
\end{tabular}
\end{center}
All random seeds are fixed, so the figures are reproducible.
